\chapter{Literature Review}
\label{ch:lit}

\section{Federated Learning}
FedAvg provides the baseline for comparing aggregation rules in this project.
In the original algorithm, participating devices start from the same model and run
local stochastic gradient descent. The server combines their resulting parameters
in proportion to the number of examples at each device
\cite{mcmahan2017fedavg}.

A sample-count weight says how much data a device holds. It says little about which
classes the device sees. Hsu et al.~\cite{hsu2019noniid} varied class imbalance
with Dirichlet sampling and observed lower FedAvg accuracy on CIFAR-10 as client
label distributions became more dissimilar. In our simulator, each class is
distributed across the UAVs using a Dirichlet draw with $\alpha=0.5$.

Abad et al.~\cite{abad2020hfl} separate aggregation within each cell from
network-wide synchronization. Their small-cell base stations combine updates from
their own devices and periodically coordinate through the macro base station.
In this project, ground base stations serve as the cell aggregators, with one
swarm-wide synchronization for every five local rounds in the controlled study.

\section{Byzantine-Robust Aggregation}
A single malicious client can move the FedAvg average arbitrarily far, so a large
body of work replaces the mean by a robust statistic.
\begin{itemize}
\item \textbf{Krum} \cite{blanchard2017krum} selects the update whose summed distance
  to its nearest neighbours is smallest. It requires $n\ge 2f+3$ clients to tolerate
  $f$ attackers.
\item \textbf{Coordinate-wise median and trimmed mean} \cite{yin2018median} take the
  median, or a trimmed mean, of each parameter separately, and come with statistical
  error guarantees.
\item \textbf{FLTrust} \cite{cao2021fltrust} gives the server a small clean root
  dataset. The server computes its own update on that data each round, scores each
  client by the (clipped at zero) cosine similarity to it, and rescales every client
  update to the norm of the server update before averaging.
\item \textbf{Geometric median} (RFA) \cite{pillutla2022rfa} aggregates with an
  approximate geometric median, which differs from the coordinate-wise median.
\end{itemize}
Non-IID data make honest updates look different from each other, which a robust rule
can mistake for an attack; bucketing \cite{karimireddy2022bucketing} adapts robust
rules to heterogeneous data. Attacks that accumulate small changes over rounds are
countered by using history, through worker momentum and iterative clipping
\cite{karimireddy2021history}; this is related to, but different from, keeping a
reputation per client.
These rules are evaluated against attacks of increasing strength. A sign-flip attack
sends a scaled negative of an honest update. ALIE (``A Little Is Enough'')
\cite{baruch2019alie} shifts each coordinate by a small multiple of the honest
standard deviation, which is enough to evade distance-based rules. Inner-product
manipulation (IPM) \cite{xie2019ipm} sends a negative multiple of the honest mean so
that the aggregate's inner product with the true gradient becomes negative.

Because the same UAVs participate across rounds in our simulator, a base station
can use their earlier updates when assigning trust. This history also gives an
adaptive adversary a way to manage its reputation. Liu et al.~\cite{liu2026evade}
formulate the decision to send a benign or malicious update as a partially
observable Markov decision process. In their experiments, an attacker able to
vary its behaviour across rounds evaded reputation-based defences and reduced
model accuracy. The controlled study in this report uses fixed attack settings.

\section{Aggregation with Wireless Erasures}
Shirvanimoghaddam et al.~\cite{shirvanimoghaddam2022erroneous}, the authors of
our base paper, ask what the server should aggregate when a device's upload is
missing. They use the finite-blocklength normal approximation to set the probability
of losing the entire upload. The server can omit that device for the round or fill
in its most recently received parameters. The authors prove convergence to the
error-free solution for reuse-old under conditions on the loss functions,
step size and erasure probability. Their uniform-data example favours the rule
that uses only fresh parameters.

Our base-paper comparison adapts both rules to uploads split across several
packets (Section~\ref{sec:method-baselines}). The server can then receive only
part of a model update. We also include Byzantine clients, whereas the original
analysis considers honest clients. Barkatsa et al.~\cite{barkatsa2024jamming} consider malicious clients that jam the
non-orthogonal multiple access (NOMA) uplink of legitimate ones, and let every client
choose its transmit power in a Bayesian game with incomplete information about the
other clients' types. Li et al.~\cite{li2022dpbyz} combine communication efficiency (1-bit
compressed sensing of sign-based gradients), differential privacy and Byzantine
robustness, but do not model lost packets. Byzantine robustness has also been studied
for analog over-the-air aggregation, in which clients transmit simultaneously and the
channel itself computes a sum: Sifaou and Li~\cite{sifaou2022ota} aggregate robustly across random
groups, Yao et al.~\cite{yao2025fedsac} combine Byzantine identification, adaptive clustering
and aggregation weights, and Fan et al.~\cite{fan2026uavbyz} use UAVs to collect gradients over
the air from heterogeneous workers under Byzantine attack. In those systems the server
cannot observe individual updates, whereas our UAVs send individually observable
digital updates that may lose packets.

\section{Federated Learning for UAV Swarms}
Xiao et al.~\cite{xiao2021uavswarm} propose fully decentralized FL for
on-board missions of a UAV swarm: UAVs exchange models with neighbours over a
time-varying graph instead of through a server, using an inexact stochastic parallel
random-walk ADMM, and the framework is demonstrated on beamforming design. It
considers neither adversarial UAVs nor packet loss. Liao et al.~\cite{liao2023tee} propose
UAV-assisted FL in which the UAV aggregates inside a trusted execution environment
(TEE), using a robust rule they call CosAvg: each client is scored by the sum of its
pairwise cosine similarities with the others, and the top-ranked updates are averaged.
Our preliminary study used a simpler CosAvg-style baseline that weights each update by
its cosine similarity to the mean update; it is not an exact reproduction. Realistic UAV link data have recently become
available: the DeepSense~6G dataset \cite{alkhateeb2023deepsense} includes a
scenario with a flying UAV and a ground base station \cite{charan2025drone}, which we
use to drive packet losses in simulation.

\section{Research Gap}
Table~\ref{tab:gap} summarizes the reviewed work; it is a small sample and does not
establish that no other work covers these points. Prior work addresses Byzantine
robustness in conventional FL, where robust rules were defined for complete update
vectors, and in analog over-the-air aggregation, including UAV-assisted
architectures. Work on erroneous links, such as the base paper, assumes honest
clients. This project studies a different setting: digital packet erasures in
hierarchical FL with UAV clients, together with model-poisoning attackers. It is also
natural to assume that link telemetry (signal quality, delivery ratio, handover state)
should make a robust aggregator \emph{fairer} to honest UAVs on poor links. This
project tests that assumption with a paired ablation that separates
delivered-coordinate masking from telemetry-weighted reputation.

\begin{table}[ht]
\centering
\caption[Coverage of the problem in the reviewed work]{Coverage of the problem in the
reviewed work. ``Poisoned updates'' means model- or data-poisoning clients;
``digital erasures'' means individually received updates that can be lost in whole or
in part.}
\label{tab:gap}
\small
\begin{tabularx}{\textwidth}{@{}>{\raggedright\arraybackslash}X
    *{4}{>{\centering\arraybackslash}p{17mm}}@{}}
\toprule
Work & Poisoned updates & Digital erasures & UAV setting & Measured channel \\
\midrule
FedAvg \cite{mcmahan2017fedavg} & -- & -- & -- & -- \\
Krum, median \cite{blanchard2017krum,yin2018median} & \checkmark & -- & -- & -- \\
FLTrust \cite{cao2021fltrust} & \checkmark & -- & -- & -- \\
Erroneous links \cite{shirvanimoghaddam2022erroneous} & -- & \checkmark & -- & -- \\
Jamming mitigation \cite{barkatsa2024jamming} & --\textsuperscript{a} & \checkmark & -- & -- \\
DP + Byzantine \cite{li2022dpbyz} & \checkmark & -- & -- & -- \\
Over-the-air robust FL \cite{sifaou2022ota,yao2025fedsac} & \checkmark & --\textsuperscript{b} & -- & -- \\
UAV-assisted over-the-air FL \cite{fan2026uavbyz} & \checkmark & --\textsuperscript{b} & \checkmark & -- \\
UAV-swarm FL \cite{xiao2021uavswarm} & -- & -- & \checkmark & -- \\
TEE + CosAvg \cite{liao2023tee} & \checkmark & -- & \checkmark & -- \\
\textbf{This project} & \checkmark & \checkmark & \checkmark & \checkmark\textsuperscript{c} \\
\bottomrule
\end{tabularx}

\raggedright\footnotesize
\textsuperscript{a}Malicious clients jam the uplink; they do not send poisoned updates.
\textsuperscript{b}Analog channel impairment of a superposed sum, not digital erasures.
\textsuperscript{c}Relative power variation of one measured link, with virtual
neighbouring BSs and assumed SNR anchors.
\end{table}
